\documentclass[10pt]{article}
\usepackage[T1]{fontenc}
\usepackage[utf8]{inputenc}
\usepackage[french]{babel}
\usepackage{amsmath,amssymb,mathrsfs,mathpazo}
\usepackage{geometry}
\geometry{a4paper,margin=2cm}
\usepackage{graphicx,array,tabularx,tabularray,colortbl,multirow,multicol,booktabs,enumitem,pifont,ulem,fancybox,diagbox,cancel}
\usepackage{pstricks,pstricks-add,pst-plot,pst-tree,pst-node,pst-text,pst-func,pst-math,pst-bspline,pst-3dplot,pst-eucl,pst-solides3d}
\let\pscalculate\undefined
\usepackage{pst-bezier}
\usepackage{tikz,pgfplots,tkz-tab}
\usetikzlibrary{arrows,arrows.meta,patterns,calc,positioning,trees,decorations.markings,shapes.geometric,shapes.misc,shapes.arrows,intersections,decorations.pathmorphing,decorations.pathreplacing,backgrounds,shadows,babel}
\pgfplotsset{compat=1.18}
\usepgfplotslibrary{fillbetween}
\usepackage[np]{numprint}
\usepackage{listings,makecell,pas-tableur,twoopt}
\let\barre\undefined
\usepackage[pstricks]{tablvar}
\definecolor{darkgreen}{RGB}{0,100,0}
\providecommand{\red}{\color{red}}
\providecommand{\blue}{\color{blue}}
\providecommand{\green}{\color{green}}
\providecommand{\black}{\color{black}}
\providecommand{\white}{\color{white}}
\providecommand{\gray}{\color{gray}}
\providecommand{\cyan}{\color{cyan}}
\providecommand{\magenta}{\color{magenta}}
\providecommand{\yellow}{\color{yellow}}
\providecommand{\violet}{\color{violet}}
\providecommand{\orange}{\color{orange}}
\providecommand{\euro}{\texteuro}
\providecommand{\EUR}{\texteuro}
\providecommand{\ieme}{\textsuperscript{e}}
\providecommand{\degres}{\ensuremath{^\circ}}
\providecommand{\degree}{\ensuremath{^\circ}}
\providecommand{\vv}[1]{\overrightarrow{#1}}
\providecommand{\Coord}[1]{#1}
\providecommand{\decosix}{}
\providecommand{\starredbullet}{\ensuremath{\bullet}}
\pagestyle{empty}
\setlength{\parindent}{0pt}
\usepackage{adjustbox}\DeclareUnicodeCharacter{2236}{:}\DeclareUnicodeCharacter{2190}{\ensuremath{\leftarrow}}
\providecommand{\R}{\mathbb{R}}
\providecommand{\N}{\mathbb{N}}
\providecommand{\Z}{\mathbb{Z}}
\providecommand{\Q}{\mathbb{Q}}
\providecommand{\D}{\mathbb{D}}
\providecommand{\C}{\mathbb{C}}
\providecommand{\e}{{\mathrm{e}}}
\providecommand{\vect}[1]{\overrightarrow{#1}}
\providecommand{\vectt}[1]{\overrightarrow{\mathrm{#1}}}
\providecommand{\vvt}[1]{\overrightarrow{\mathrm{#1}}}
\providecommand{\np}[2][]{{#2~#1}}
\providecommand{\vv}[1]{{\overrightarrow{#1}}}
\renewcommand{\d}{{\mathrm{d}}}
\providecommand{\degres}{{^\circ}}
\providecommand{\degre}{{^\circ}}
\providecommand{\blue}{\color{blue}}
\providecommand{\red}{\color{red}}
\providecommand{\enskip}{\quad}
\providecommand{\llbracket}{\lbrack\!\lbrack}
\providecommand{\rrbracket}{\rbrack\!\rbrack}
\providecommand{\euro}{{\text{€}}}
\providecommand{\no}{{\text{n°}}}
\providecommand{\interval}[2]{{[#1;#2]}}
\providecommand{\intervFO}[2]{{[#1;#2[}}
\providecommand{\dotfill}{\ldots}

\providecommand{\ds}{\displaystyle}
\providecommand{\textding}[1]{#1}
\providecommand{\Oij}{\(\mathrm{O};\vec{\imath},\vec{\jmath}\)}
\providecommand{\Oijk}{\(\mathrm{O};\vec{\imath},\vec{\jmath},\vec{k}\)}
\providecommand{\Ouv}{\(\mathrm{O};\vec{u},\vec{v}\)}
\providecommand{\intervFF}[2]{\ensuremath{[#1;#2]}}
\renewcommand{\euro}{\ifmmode\text{\texteuro}\else\texteuro\fi}
\DeclareUnicodeCharacter{00B0}{\ifmmode^{\circ}\else\textdegree\fi}
\DeclareUnicodeCharacter{20AC}{\ifmmode\text{\texteuro}\else\texteuro\fi}
% Énoncés : Éducation nationale. Transcription et corrigés : APMEP, Denis Vergès et contributeurs. Découpage et classement : CoopMaths, Sébastien Lozano, Éric Elter et contributeurs. Adaptation de la présentation : Maths Explorer.
% CC BY-SA : https://coopmaths.fr/www/about/
% https://coopmaths.fr/alea/static/e3c/2020/tex/e3c_2020_00_sujet55_3_cor.tex
\begin{document}
\noindent\textit{APMEP / CoopMaths. CC BY-SA. Présentation adaptée par Maths Explorer.}\par\medskip
\subsection*{1.}

\begin{align*}
&M(x\,;\,y) \in d_1 \\
\iff &\overrightarrow{CM} \cdot \overrightarrow{AB} = 0 \\
\iff &0 \times (x - 1) + 6(y - 1) = 0 \\
\iff &y - 1 = 0.
\end{align*}

\subsection*{2.}

La droite contenant un sommet et perpendiculaire au côté opposé est une hauteur.

\subsection*{3.}

\begin{align*}
&M(x\,;\,y) \in d_2 \\
\iff &\overrightarrow{BM} \cdot \overrightarrow{AC} = 0 \\
\iff &-6(x - 7) + 3(y - 4) = 0 \\
\iff &-6x + 3y + 42 - 12 = 0 \\
\iff &-2x + y + 10 = 0.
\end{align*}

\subsection*{4.}

Si \( H \) est le point d'intersection de deux hauteurs, c'est donc l'orthocentre, point commun aux trois hauteurs. La troisième hauteur est donc \( (AH) \), qui est perpendiculaire à \( (BC) \). Les vecteurs \( \overrightarrow{AH} \) et \( \overrightarrow{BC} \) sont donc orthogonaux, et leur produit scalaire est donc nul.

\begin{center}
\psset{xunit=0.5cm,yunit=0.5cm,labelsep=0.1pt,labelFontSize=\scriptstyle,showorigin=false}
\begin{pspicture}(-2,-4.5)(11,8)
 \multido{\n=-1+1}{13}{\psline[linewidth=0.75pt,linecolor=lightgray](\n,-4)(\n,7.5)}
 \multido{\n=-4+1}{12}{\psline[linewidth=0.75pt,linecolor=lightgray](-1.5,\n)(11.2,\n)}
 \psaxes[linewidth=0.95pt]{->}(0,0)(-1.5,-4)(11,7.5)
\uput[dl](0,-0.2){\footnotesize O}\uput[l](6.8,-2.5){A}\uput[l](6.8,4.5){B}\uput[l](0.8,0.5){C}\uput[l](5.2,1.4){H}
\psdots[dotstyle=+,dotscale =1.4,dotangle=45](1,1)(7,4)(7,-2)
\psline[linewidth=1pt](7,-2)(7,4)
\psline[linewidth=1pt](7,4)(1,1)
\psline[linewidth=1pt](1,1)(7,-2)
\psline[linewidth=1pt,linecolor=blue](1,1)(8,1)
\psline[linewidth=1pt,linecolor=red](8,-4)(4,4)
\psline[linewidth=1pt,linecolor=blue](4,-2)(8,6)
\end{pspicture}
\end{center}

\end{document}
